\section{Introduccion}

\subsection{Caso de Estudio}
La defensa cibernética es un conjunto de estrategias, tecnologías y procesos
diseñados para proteger los sistemas informáticos, redes y datos contra
amenazas cibernéticas, como ataques de malware, piratería informática, robo de
datos y otros riesgos de seguridad. Estas medidas de defensa buscan prevenir,
detectar, responder y recuperarse de ataques cibernéticos con el fin de
garantizar la confidencialidad, integridad y disponibilidad de la información y
los recursos digitales.

\subsection{Resumen Ejecutivo}

En el contexto actual de la transformación digital, la ciberseguridad se ha convertido en un pilar fundamental para la protección de los sistemas de información. Las amenazas informáticas evolucionan constantemente, lo que exige el desarrollo de soluciones inteligentes capaces de detectar, analizar y responder a incidentes de seguridad de manera eficiente.

Este proyecto presenta el desarrollo inicial de un \textbf{agente inteligente de ciberseguridad}, implementado utilizando tecnologías modernas como \textbf{Python}, el framework web \textbf{Django} y el sistema de gestión de bases de datos \textbf{PostgreSQL}. Aunque el agente aún se encuentra en una etapa temprana y su funcionalidad es limitada, se ha logrado establecer una arquitectura básica que permite la integración y gestión de datos relevantes para la detección de amenazas.

El sistema está diseñado para operar como un agente autónomo dentro de una infraestructura web, recolectando información sobre eventos y posibles vulnerabilidades. La elección de Django facilita la construcción de una interfaz administrativa y la gestión eficiente de los modelos de datos, mientras que PostgreSQL proporciona robustez y escalabilidad en el almacenamiento de información.

Cabe resaltar que, si bien el agente aún no implementa técnicas avanzadas de inteligencia artificial, se han sentado las bases para futuras mejoras, incluyendo la incorporación de algoritmos de aprendizaje automático y módulos de análisis automatizado. El objetivo principal de este informe es documentar el proceso de diseño y desarrollo, así como los retos enfrentados y las oportunidades de mejora identificadas en el camino.
